


\usepackage{amssymb}
%\usepackage[T1]{fontenc}
%\usepackage[utf8]{inputenc}

%\usepackage[dvipsnames]{xcolor}

%\usepackage{pifont}
%\usepackage{todonotes}
\usepackage{mathtools}
\usepackage{stmaryrd}
\usepackage{xspace}
\usepackage{hyperref}
\usepackage[nameinlink]{cleveref}

%\usepackage{fdsymbol}

\usepackage[shortcuts]{extdash}

\usepackage{tikz}


\hyphenation{Skrzyp-czak}

% QED

% Operator

\newcommand{\ident}[1]{\ensuremath{\texorpdfstring{\mathrm{\mathnm{#1}}}{#1}}\xspace}
\newcommand{\operator}[1]{\ensuremath{\texorpdfstring{\mathop{#1}\xspace}{#1}}}

% Domain

\newcommand{\domain}[1]{\ensuremath{\mathbb{#1}}\xspace}

\newcommand{\sem}[1] {  \llbracket {#1} \rrbracket  }  

\newcommand{\ot}{\leftarrow}

\newcommand{\lex}{<_{\mathrm{\mathnm{lex}}}}
\newcommand{\lexeq}{\leq_{\mathrm{\mathnm{lex}}}}
\newcommand{\lexgeq}{\geq_{\mathrm{\mathnm{lex}}}}
\newcommand{\lexge}{>_{\mathrm{\mathnm{lex}}}}

\newcommand{\wadge}{<_{\mathrm{\mathnm{W}}}}
\newcommand{\wadgeq}{\leq_{\mathrm{\mathnm{W}}}}
\newcommand{\wadggeq}{\geq_{\mathrm{\mathnm{W}}}}
\newcommand{\wadgesim}{\equiv_{\mathrm{\mathnm{W}}}}

\newcommand{\KBe}{<_{\mathrm{\mathnm{KB}}}}
\newcommand{\KBeq}{\leq_{\mathrm{\mathnm{KB}}}}
\newcommand{\KBge}{>_{\mathrm{\mathnm{KB}}}}
\newcommand{\KBgeq}{\geq_{\mathrm{\mathnm{KB}}}}


\newcommand{\DOT}{$\bullet\ $}
\newcommand{\ARR}{{\Large${\pmb{\rightarrow}}$}\xspace}
\newcommand{\RRA}{{\Large${\pmb{\leftarrow}}$}\xspace}

\newcommand{\WG}{\hl{$\pmb{\rightsquigarrow}$}\xspace}
\newcommand{\GW}{\hl{$\pmb{\leftsquigarrow}$}\xspace}


\newcommand{\qU}{\ensuremath{\mathsf U}\xspace}

\newcommand{\qG}{\rotatebox[origin=c]{180}{G}\xspace}

\newcommand{\concat}{\text{\^{}}}

\newcommand{\mysc}[1]{{\sc {\texorpdfstring{\mathnm{#1}}{#1}}}\xspace}

\newcommand{\zfc}{\mysc{zfc}}
\newcommand{\fo}{\mysc{fo}}
\newcommand{\rca}{\ensuremath{\text{\mysc{rca}}_0}}
\newcommand{\ef}{\mysc{ef}}
\newcommand{\emso}{\ensuremath{\exists}\mysc{mso}}
\newcommand{\mso}{\mysc{mso}}
\newcommand{\msoR}{\mysc{mso$(\R)$}}
\newcommand{\ltl}{\mysc{ltl}}
\newcommand{\ctl}{\mysc{ctl}}
\newcommand{\ctlS}{\mysc{ctl*}}
\newcommand{\msou}{\mysc{mso+u}}
\newcommand{\wmso}{\mysc{wmso}}
\newcommand{\wmsou}{\mysc{wmso+u}}
\newcommand{\wmsor}{\mysc{wmso+r}}
\newcommand{\wmsoup}{\mysc{wmso+up}}
\newcommand{\VL}{\mysc{v=l}}

\newcommand{\NL}{\mysc{NL}}
\newcommand{\SL}{\mysc{SL}}
\newcommand{\LOGSPACE}{\mysc{LogSpace}}
\newcommand{\NP}{\mysc{NP}}
\newcommand{\UP}{\mysc{UP}}
\newcommand{\coNP}{\mysc{co-NP}}
\newcommand{\PTIME}{\mysc{PTime}}
\newcommand{\Pt}{\mysc{P}}
\newcommand{\PSPACE}{\mysc{PSpace}}
\newcommand{\NPSPACE}{\mysc{NPSpace}}
\newcommand{\SPACE}{\mysc{Space}}
\newcommand{\NSPACE}{\mysc{NSpace}}
\newcommand{\EXPTIME}{\mysc{EXPTime}}


\newcommand{\MA}{\ensuremath{\mathbf {MA}_{\aleph_1}}\xspace}

\newcommand{\projmso}{\mbox{proj-\mysc{mso}}\xspace}

\newcommand{\powerset}{\mathsf{\mathnm{P}}}

% w-BS
\newcommand{\wBS}{\ensuremath{\omega \mathrm{BS}}\xspace}
\newcommand{\wB}{\ensuremath{\omega \mathrm{B}}\xspace}
\newcommand{\wS}{\ensuremath{\omega \mathrm{S}}\xspace}

\newcommand{\wT}{\ensuremath{\omega \mathrm{T}}\xspace}

\newcommand{\fT}{\ensuremath{\mathrm{T}}\xspace}
\newcommand{\fB}{\ensuremath{\mathrm{B}}\xspace}
\newcommand{\fS}{\ensuremath{\mathrm{S}}\xspace}

% Numbers

\newcommand{\A}{\domain{A}}
\newcommand{\B}{\domain{B}}
\newcommand{\C}{\domain{C}}
\newcommand{\D}{\domain{D}}
\newcommand{\E}{\domain{E}}
\newcommand{\N}{\domain{N}}
\newcommand{\Z}{\domain{Z}}
\newcommand{\Q}{\domain{Q}}
\newcommand{\R}{\domain{R}}
\newcommand{\Pri}{\domain{P}}

%\DeclareMathOperator*{\prob}{\mathbb P}

\newcommand{\prom}[1]{\ensuremath{\widehat{{#1}^\ast}}\xspace}

\newcommand{\aut}[1]{\ensuremath{\mathcal {#1}}\xspace}

% Borel and analytic sets

\newcommand{\boldclass}[3]{\ensuremath{\mathbf{#1}^{#2}_{#3}}}
\newcommand{\lightclass}[3]{\ensuremath{{#1}^{#2}_{#3}}}

\newcommand{\borel}{\ensuremath{\mathcal B}\xspace}

\newcommand{\bsigma}[1]{\boldclass{\Sigma}{0}{#1}}
\newcommand{\bpi}[1]{\boldclass{\Pi}{0}{#1}}
\newcommand{\bdelta}[1]{\boldclass{\Delta}{0}{#1}}

\newcommand{\asigma}[1]{\boldclass{\Sigma}{1}{#1}}
\newcommand{\api}[1]{\boldclass{\Pi}{1}{#1}}
\newcommand{\adelta}[1]{\boldclass{\Delta}{1}{#1}}

\newcommand{\esigma}[2]{\lightclass{\Sigma}{#1}{#2}}
\newcommand{\epi}[2]{\lightclass{\Pi}{1}{#1}}
\newcommand{\edelta}[2]{\lightclass{\Delta}{#1}{#2}}

\newcommand{\bc}[1]{\mathcal{BC}({#1})}

\newcommand{\RM}{\ident{\mathbf{RM}^{\mathrm{alt}}}}

\newcommand{\RMnd}{\ident{\mathbf{RM}^{\mathrm{non-det}}}}

\newcommand{\RMeven}[1]{\ensuremath{\mathbf{\Pi}^{\mathrm{alt}}_{#1}}\xspace}
\newcommand{\RModd}[1]{\ensuremath{\mathbf{\Sigma}^{\mathrm{alt}}_{#1}}\xspace}
\newcommand{\RMdelta}[1]{\ensuremath{\mathbf{\Delta}^{\mathrm{alt}}_{#1}}\xspace}
\newcommand{\RMcmp}[1]{\ensuremath{\mathbf{Comp}(\RMeven{#1})}\xspace}

\newcommand{\cmp}{\ident{Comp}}
\newcommand{\graph}{\ident{Graph}}

% def

\newcommand{\eqdef}{\stackrel{\text{def}}=}
\newcommand{\eqqu}{\stackrel{\text{?}}=}

\newcommand{\eqidef}{\stackrel{\text{def}}\equiv}


\newcommand{\Efin}{\exists^{\mathrm{fin}}}
\newcommand{\Afin}{\forall^{\mathrm{fin}}}


\renewcommand{\mod}[1]{\allowbreak\mkern10mu({\operator@font mod}\,\,#1)}

\newcommand*{\plainfrac}[2]{\genfrac{}{}{0pt}{}{#1}{#2}}

% Roman numbers

\newcommand{\rI}{\mathrm{I}}
\newcommand{\rII}{\mathrm{II}}
\newcommand{\rIII}{\mathrm{III}}
\newcommand{\rIV}{\mathrm{IV}}
\newcommand{\rV}{\mathrm{V}}

% numbers

\newcommand{\mathcalsym}[1]{\ensuremath{\mathcal{#1}}\xspace}

\newcommand{\Chi}{\scalerel*{\chi}{\sum}}

\newcommand{\CA}{\mathcalsym{A}}
\newcommand{\CB}{\mathcalsym{B}}
\newcommand{\CC}{\mathcalsym{C}}
\newcommand{\CD}{\mathcalsym{D}}
\newcommand{\CE}{\mathcalsym{E}}
\newcommand{\CF}{\mathcalsym{F}}
\newcommand{\CG}{\mathcalsym{G}}
\newcommand{\CH}{\mathcalsym{H}}
\newcommand{\CI}{\mathcalsym{I}}
\newcommand{\CJ}{\mathcalsym{J}}
\newcommand{\CK}{\mathcalsym{K}}
\newcommand{\CL}{\mathcalsym{L}}
\newcommand{\CM}{\mathcalsym{M}}
\newcommand{\CN}{\mathcalsym{N}}
\newcommand{\CO}{\mathcalsym{O}}
\newcommand{\CP}{\mathcalsym{P}}
\newcommand{\CQ}{\mathcalsym{Q}}
\newcommand{\CR}{\mathcalsym{R}}
\newcommand{\CS}{\mathcalsym{S}}
\newcommand{\CT}{\mathcalsym{T}}
\newcommand{\CU}{\mathcalsym{U}}
\newcommand{\CV}{\mathcalsym{V}}
\newcommand{\CW}{\mathcalsym{W}}
\newcommand{\CX}{\mathcalsym{X}}
\newcommand{\CY}{\mathcalsym{Y}}
\newcommand{\CZ}{\mathcalsym{Z}}

% Greek letters

\newcommand{\Ga}{\alpha}
\newcommand{\Gb}{\beta}
\newcommand{\Gg}{\gamma}
\newcommand{\Gd}{\delta}
\newcommand{\Ge}{\epsilon}
\newcommand{\Gve}{\varepsilon}
\newcommand{\Gz}{\zeta}
\newcommand{\Gn}{\eta}
\newcommand{\Gm}{\mu}
\newcommand{\Gh}{\theta}
\newcommand{\Gk}{\kappa}
\newcommand{\Gl}{\lambda}
\newcommand{\Gu}{\nu}
\newcommand{\Gx}{\xi}
\newcommand{\Gp}{\pi}
\newcommand{\Gy}{\upsilon}
\newcommand{\Gt}{\tau}
\newcommand{\Gs}{\sigma}
\newcommand{\Gr}{\rho}
\newcommand{\Gw}{\omega}

% Tikz

% vertical subseteq
\newcommand{\vsubseteq}{\rotatebox[origin=c]{90}{$\subseteq$}}
\newcommand{\vsupseteq}{\rotatebox[origin=c]{90}{$\supseteq$}}
\newcommand{\vequiv}{\rotatebox[origin=c]{90}{$\equiv$}}
\newcommand{\vcong}{\rotatebox[origin=c]{90}{$\cong$}}
\newcommand{\vpreceq}{\rotatebox[origin=c]{90}{$\preceq$}}
\newcommand{\vprec}{\rotatebox[origin=c]{90}{$\prec$}}

\renewcommand{\boldsymbol}{\pmb}

% Functions


\newcommand{\from}{\colon}

\newcommand{\fun}[3]{\ensuremath{#1\from #2 \to #3}}
\newcommand{\funmap}[3]{\ensuremath{#1\from #2 \mapsto #3}}
\newcommand{\parfun}[3]{\ensuremath{#1\from #2 \rightharpoonup #3}}

\newcommand{\dom}{\ident{dom}}
\newcommand{\rg}{\ident{rg}}
\newcommand{\id}{\ident{id}}
\newcommand{\sgn}{\ident{sgn}}

\newcommand{\close}{\ident{cl}}
\newcommand{\inter}{\ident{int}}
\newcommand{\lang}{\ident{L}}

% language W like Walukiewicz
\newcommand{\W}[2]{\ensuremath{W_{{#1},{#2}}}\xspace}

%\newcommand{\dL}{\mathbf{L}}
%\newcommand{\dR}{\mathbf{R}}


\usetikzlibrary{positioning}
\usetikzlibrary{decorations.pathreplacing}
\usetikzlibrary{automata,positioning}
\usetikzlibrary{decorations.pathmorphing}
\usetikzlibrary{decorations.markings}
\usetikzlibrary{decorations}
\usetikzlibrary{arrows}
\usetikzlibrary{patterns}
\usetikzlibrary{calc}
\usetikzlibrary{shapes}
\usetikzlibrary{fadings,	shadings}

\usetikzlibrary{arrows.meta}
%\usetikzlibrary{snakes}

	\tikzstyle{ubrace} = [draw, thick, decoration={brace, amplitude=7pt, mirror, raise=0.0cm}, decorate,
	    every node/.style={anchor=north, yshift=-0.15cm}]
	\tikzstyle{rbrace} = [draw, thick, decoration={brace, amplitude=7pt, mirror, raise=0.0cm}, decorate,
	    every node/.style={anchor=west, xshift= 0.15cm}]

	\tikzstyle{obrace} = [draw, thick, decoration={brace, amplitude=7pt, raise=0.0cm}, decorate,
	    every node/.style={anchor=south, yshift= 0.15cm}]
	\tikzstyle{lbrace} = [draw, thick, decoration={brace, amplitude=7pt, raise=0.0cm}, decorate,
	    every node/.style={anchor=east, xshift=-0.15cm}]



\newcommand{\eval}[2]{%
	\pgfmathparse{#2}%
	{\global\edef#1{\pgfmathresult}}%
}

% This one returns the number rounded as an integer
\newcommand{\evalInt}[2]{%
	\pgfmathparse{int(#2)}%
	{\global\edef#1{\pgfmathresult}}%
}


\newcommand{\dL}{\mathtt{{\scriptstyle L}}}
\newcommand{\dR}{\mathtt{{\scriptstyle R}}}

\newcommand{\mathnm}[1]{#1}

\newcommand{\hole}{\Box}

%\newcommand{\tran}[1]{\overset{#1}{\longrightarrow}}
\newcommand{\tran}[1]{\xrightarrow{#1}}

% complement
\newcommand{\cp}{\mathrm{\mathnm{c}}}

% Operations

\newcommand{\restr}{{\upharpoonright}}

\newcommand{\newton}[2]{\ensuremath{\binom{#1}{#2}}}

\newcommand{\clos}[1]{\ensuremath{\overline{#1}}}

\newcommand{\implication}[2]{\texorpdfstring
{$\eqref{#1}\Rightarrow\eqref{#2}$}
{(\ref{#1})=>(\ref{#2})}}

\tikzstyle{dot} = [draw,shape=circle,fill, minimum size=1mm, inner sep=0pt, outer sep=0pt]
\tikzstyle{edge}  = [draw, thick,->]




\newenvironment{proofof}[1]
	{\vspace{1ex}\noindent{\emph{Proof of #1}}\hspace{0.5em}}
    {\hfill\qed\vspace{1ex}}
    




\newcommand{\init}{\ident{I}}

\DeclareSymbolFont{yhlargesymbols}{OMX}{yhex}{m}{n}
\DeclareMathAccent{\yhwidehat}{\mathord}{yhlargesymbols}{"62}

\renewcommand{\widehat}{\yhwidehat}

\newcommand{\trees}{\ident{Tr}}
\newcommand{\hist}{\ident{hist}}
\newcommand{\val}{\ident{val}}
\newcommand{\oal}{\ident{val}^{+1}}
\newcommand{\Pf}{\mathit{Pf}}

\newcommand{\WF}{\ident{WF}}

\newcommand{\der}{\ident{der}}

\newcommand{\unravel}{\ident{unravel}}

\newcommand{\sT}{\mathsf{T}}
\newcommand{\sR}{\mathsf{R}}

\newcommand{\rank}{\ident{rank}}
\newcommand{\rankX}[1]{\ident{rank}^{#1}}
\newcommand{\rankR}{\ident{rank}^\sR}
\newcommand{\rankT}{\ident{rank}^\sT}
\newcommand{\apr}{\ident{apr}}

\renewcommand\phi\varphi
\renewcommand\epsilon\varepsilon
\newcommand\usun[1]{\begingroup\color{red}{}#1\endgroup}

\newcommand\zmiana[1]{\begingroup\color{blue}{}#1\endgroup}

\newcommand{\eve}{\ensuremath{\exists}\xspace}
\newcommand{\adam}{\ensuremath{\forall}\xspace}



\newcommand{\ignore}[1]{}

\newcommand*\circled[1]{\tikz[baseline=(char.base),scale=0.6]{
            \node[shape=circle,draw,inner sep=0mm,scale=0.83,minimum size=3.7mm] (char) {\textbf{#1}};}}


\newcommand{\fb}{\mathbf{\text{\circled{2}}}}
\newcommand{\fg}{\mathbf{\text{\circled{1}}}}



\DeclarePairedDelimiter{\set}{\{}{\}}
\DeclarePairedDelimiter{\pair}{\llparenthesis}{\rrparenthesis}

\newcommand{\imwin}[2]{[#1\to #2]}

%\newcommand\dn[1]{\todo[size=\scriptsize]{#1 - \textbf{Damian}}}
%\newcommand\dn[1]{\todo[size=\scriptsize,backgroundcolor=Green]{#1 -\textbf{Damian}}}
\newcommand\dn[1]{\todo[inline,size=\scriptsize,backgroundcolor=green]{#1 -\textbf{Damian}}}
\newcommand\pp[1]{\todo[inline, size=\scriptsize,backgroundcolor=yellow]{#1 - \textbf{Paweł}}}
\newcommand\ms[1]{\todo[inline, size=\scriptsize,backgroundcolor=cyan]{#1 - \textbf{Michał}}}

\Crefname{lem}{Lemma}{Lemmata}
\Crefname{cor}{Corollary}{Corollaries}
\Crefname{thm}{Theorem}{Theorems}
\Crefname{fact}{Fact}{Facts}
\Crefname{clm}{Claim}{Claims}
\Crefname{rem}{Remark}{Remarks}
\Crefname{exa}{Example}{Examples}
\Crefname{qu}{Question}{Questions}
\Crefname{prop}{Proposition}{Propositions}
\Crefname{equation}{Formula}{Formulae}
\Crefname{property}{Property}{Properties}\creflabelformat{property}{(#2#1#3)}

\newcommand{\bs}[1]{\mathbf{#1}}

\newcommand{\dist}{\mathit{dist}}
\newcommand{\stab}{\mathit{st}}

\newcommand{\wit}{\ident{wit}_\omega}

\newcommand{\limit}{\ident{rank}}


% \newcommand{\derR}{\der_1}
\newcommand{\derF}{\der_\omega}

\newcommand{\limitR}{\mathrm{rank}}
\newcommand{\limitF}{\mathrm{rank}^{\sR}}




% ---- TIKZ ----

\tikzstyle{very loosely dotted}=[dash pattern=on \pgflinewidth off 6pt]

\newcommand{\gSta}[2]{
	\node[anchor=mid] at #1 {#2};
}

\newcommand{\gStT}[4]{
	\gSta{(#1*\ssx,#2*\ssy+#3*0.35)}{#4}
}

\newcommand{\gStR}[4]{
	\gSta{(#1*\ssx+\sss,#2*\ssy+#3*0.35)}{#4}
}

\newcommand{\gDot}[2]{
	\node[dot] (#1) at #2 {};
}


\newcommand{\gCon}[2]{
	\draw (#1) -- (#2);
}

\newcommand{\gConP}[2]{
	\draw (#1) edge[dashed] (#2);
}

\newcommand{\gLab}[3]{
	\node[anchor=mid east] at (-\ssx, #3*\ssy+0.5*\ssy) {#1:};
	\node[anchor=mid west] at (-\ssx, #3*\ssy+0.5*\ssy) {#2};
}

\newcommand{\gTta}[4]{
	\gDot{#1-T}{($(#1)+(+#2*0.3*\ssx,+1*\ssy)$)}
	\gDot{#1-R}{($(#1)+(-#2*0.3*\ssx,+1*\ssy)$)}

	\node[scale=0.8, anchor=mid east] at ($(#1)+(-0.2,-0.2)$) {$\delta_{#4}$};
	\node[scale=0.8, anchor=mid] at ($(#1)+(0,+0.95*\ssy)$) {#3};
}

\newcommand{\gTtb}[4]{
	\gTta{#1}{#2}{#3}{#4}
	
	\draw (#1) edge[line width=0.6mm, line cap=round] (#1-T);
	\draw (#1) edge[] (#1-R);
}

\newcommand{\gTtl}[4]{
	\gTta{#1}{#2}{#3}{#4}
	
	\draw (#1) edge[](#1-T);
	\draw (#1) edge[very loosely dotted](#1-R);
}

\newcommand{\gTrl}[4]{
	\gTta{#1}{#2}{#3}{#4}
	
	\draw (#1) edge[](#1-T);
	\draw (#1) edge[very loosely dotted](#1-R);
}
